\documentclass[11pt]{article}
\usepackage[T1]{fontenc}
\usepackage{textcomp}
\usepackage[margin=0.75in]{geometry}
\usepackage{amsmath,amssymb,amsthm}
\usepackage{array,booktabs}
\usepackage{hyperref}
\setlength{\parindent}{0pt}
\newtheorem{theorem}{Theorem}
\theoremstyle{definition}
\newtheorem{definition}{Definition}
\title{Flow Proof Artifact: prop-26-similar-parallelograms-duplicate-ratio}
\author{Generated by \texttt{flow doc proof}}
\date{Flow Proof Book}
\begin{document}
\maketitle
Similar parallelograms are to one another in the duplicate ratio of the homologous sides.\\[0.5em]
\textbf{Source.} euclid — Elements, Book VI, Proposition 26\\[0.5em]
\bigskip
\noindent{\large \textbf{Derived fact 1.} \textit{Similar parallelograms are to one another in the duplicate ratio of the homologous sides} \hfill the Euclidean plane \cdot Euclid Book VI \cdot Proposition 26: similar parallelograms are in duplicate ratio of homologous sides}\par
\smallskip\noindent\textit{Coordinate: the Euclidean plane \cdot Euclid Book VI \cdot Proposition 26: similar parallelograms are in duplicate ratio of homologous sides}\par
\smallskip\noindent\textit{Source: euclid — Elements, Book VI, Proposition 26}\par
\smallskip\noindent\textit{Built on: proposition 19: similar rectilineal figures are in duplicate ratio of corresponding sides, for Euclid Book VI on the Euclidean plane, proposition 24: parallelograms about a diameter are similar to the whole and to one another, for Euclid Book VI on the Euclidean plane}\par
\medskip
\noindent\textbf{Goal.} Similar parallelograms are to one another in the duplicate ratio of the homologous sides\par
\noindent$\displaystyle \text{similar parallelograms duplicate ratio}$\par
\medskip
\noindent
\begin{tabular}{@{} >{\raggedright\arraybackslash}p{0.06\textwidth} >{\raggedright\arraybackslash}p{0.47\textwidth} >{\raggedleft\arraybackslash}p{0.41\textwidth} @{}}
\textbf{\#} & \textbf{Proof} & \textbf{Mathematics} \\
\midrule
\textbf{1.} & We prove that proposition 26: similar parallelograms are in duplicate ratio of homologous sides for Euclid Book VI the Euclidean plane. & \\[0.45em]
\textbf{2.} & Let par1 = parallelogram\_1. & \\[0.45em]
\textbf{3.} & Let par2 = parallelogram\_2. & \\[0.45em]
\textbf{4.} & Let side = homologous\_side. & \\[0.45em]
\textbf{5.} & From step 2, step 3, and step 4, we can deduce that area ratio duplicate. & $\displaystyle \text{area ratio duplicate}$ \\[0.45em]
\textbf{6.} & We invoke the derived fact governing Euclid Book VI the Euclidean plane: proposition 19: similar rectilineal figures are in duplicate ratio of corresponding sides, for Euclid Book VI on the Euclidean plane. & \\[0.45em]
\textbf{7.} & We invoke the derived fact governing Euclid Book VI the Euclidean plane: proposition 24: parallelograms about a diameter are similar to the whole and to one another, for Euclid Book VI on the Euclidean plane. & \\[0.45em]
\textbf{8.} & From step 6 and step 7, this implies similar parallelograms duplicate ratio. Hence proven. & $\displaystyle \text{similar parallelograms duplicate ratio}$ \\[0.45em]
\bottomrule
\end{tabular}
\label{thm:Geometry-Euclid Book VI-proposition_26_similar_parallelograms_are_in_duplicate_ratio_of_homologous_sides}
\par\medskip\hrule
\end{document}
